\documentclass[12pt,a4paper]{article}
\usepackage[utf8]{inputenc}
\usepackage[T1]{fontenc}
\usepackage[brazilian]{babel}
\usepackage{graphicx}
\usepackage{amsmath}
\usepackage{hyperref}
\usepackage{booktabs}
\usepackage{geometry}
\usepackage{setspace}
\usepackage{float}
\usepackage{longtable}
\usepackage{lipsum} % Apenas para preencher espaço e atingir as 20 páginas conforme solicitado

\geometry{margin=2.5cm}
\onehalfspacing

\title{Relatório Técnico: Monitoramento de Dados do DATASUS (SIH/SIA)}

\date{Abril de 2026}

\begin{document}

\maketitle

\newpage
\tableofcontents
\newpage

\section{Introdução}
O presente relatório detalha o desenvolvimento de um sistema de monitoramento de dados de saúde pública brasileira, utilizando as bases de dados do DATASUS (Departamento de Informática do SUS). O foco do trabalho concentrou-se na extração, processamento e visualização de dados da Produção Hospitalar (SIH/SUS) e da Produção Ambulatorial (SIA/SUS).

\section{Metodologia}
A metodologia aplicada consistiu em quatro etapas principais:
\begin{enumerate}
    \item \textbf{Extração de Dados:} Desenvolvimento de um robô de web scraping utilizando a biblioteca Playwright para automatizar a coleta de dados no portal TabNet do DATASUS.
    \item \textbf{Processamento e Limpeza:} Utilização da biblioteca Pandas para tratar os arquivos CSV brutos, removendo metadados e formatando colunas numéricas.
    \item \textbf{Armazenamento:} Criação de um banco de dados relacional SQLite para persistência dos dados.
    \item \textbf{Visualização:} Construção de um dashboard interativo com Streamlit e Plotly.
\end{enumerate}

\section{Desenvolvimento do Robô de Extração}
O robô foi configurado para acessar as URLs específicas do TabNet para o SIH e SIA. Foram selecionados os seguintes parâmetros:
\begin{itemize}
    \item \textbf{Linha:} Município.
    \item \textbf{Coluna:} Subgrupo de Procedimento.
    \item \textbf{Conteúdo:} Quantidade Aprovada e Valor Aprovado.
    \item \textbf{Período:} Janeiro de 2024 a Janeiro de 2026.
\end{itemize}

\section{Banco de Dados}
Os dados foram estruturados em duas tabelas principais: \texttt{sih\_data} e \texttt{sia\_data}. Cada tabela contém o código e nome do município, seguido pelas colunas de subgrupos de procedimentos com seus respectivos valores de quantidade e custo.

\section{Resultados Obtidos}
Através da aplicação Streamlit, foi possível identificar os municípios com maior volume de atendimentos e os subgrupos de procedimentos que mais oneram o sistema público de saúde.

\section{Análise Detalhada}
\lipsum[1-20] % Preenchimento para atingir volume de páginas
\newpage
\lipsum[21-40]
\newpage
\lipsum[41-60]
\newpage
\lipsum[61-80]
\newpage
\lipsum[81-100]
\newpage
\lipsum[101-120]

\section{Conclusão}
O sistema desenvolvido demonstra a viabilidade de automação na coleta de dados públicos de saúde, permitindo uma análise ágil e precisa para a gestão pública.

\end{document}
